\section{Introduction}
\label{sec:introduction}

Forecasting has recently been reframed as an agentic activity. Rather than fitting a
model and emitting a prediction, a growing family of systems orchestrates diagnostics,
selects among candidate models, produces a forecast, and then \emph{reflects} on it,
revising where the reflection suggests a problem
\citep{tsscientist2025,nexus2026,castflow2026,flairr2025,timeclaw2026}. A position
paper has argued explicitly that forecasting should move beyond model-centric
prediction toward iterative, tool-using workflows \citep{position2026}.

The reflection stage is now close to standard. Its value has not been established.

This is not an oversight by any individual paper so much as a structural gap in how the
area is evaluated. Three obstacles stand in the way of assessing self-correction.

\paragraph{The correction decision has no ground truth.}
Benchmarks record the error of the forecast an agent finally produced. They do not
record what the error \emph{would} have been under a different corrective action, or
under none. Without those counterfactuals, ``the agent corrected and the error was
$x$'' cannot be turned into ``correcting was the right call''. Precision and recall of
the decision to intervene, the regret of choosing one repair over another, and the rate
at which corrections actively cause harm are all unmeasurable.

\paragraph{Comparisons are confounded by computation.}
A correcting agent runs more models, calls more tools, and consumes more tokens than a
non-correcting one. When it wins, the win is jointly attributable to better decisions
and to a larger budget. Almost no published comparison in this area separates the two,
which means the headline claim---that self-correction helps---has an unexamined
alternative explanation.

\paragraph{Agents are asked to judge themselves.}
Reflection in these systems is typically verbal: the model is prompted to assess its own
output. For reasoning tasks this is precisely the mechanism shown to be unreliable.
\citet{huang2024selfcorrect} found that intrinsic self-correction, without external
feedback, often fails to improve and can degrade accuracy, and
\citet{kamoi2024selfcorrect} established that the gains which do materialise come from
outside the model---a tool, an execution result, an environment signal. Whether this
transfers to forecasting is an open and consequential question, because forecasting is
unusual: a cheap, principled external signal is always available in the form of
backtesting on recent history.

\paragraph{Classical forecasting already knows the answer to a neighbouring question.}
Long before agents, the forecasting literature established that \emph{combining}
forecasts beats trying to pick the right one. \citet{bates1969} showed that a simple
average of two forecasts can beat both; five decades of subsequent work
\citep{clemen1989,timmermann2006,wang2023review} and every large forecasting competition
\citep{makridakis2000m3,m4competition2020} have confirmed it. Sharper still is the
\emph{forecast combination puzzle}: equally-weighted combinations routinely outperform
theoretically optimal weights estimated from data
\citep{stockwatson2004,smithwallis2009,claeskens2016}, because estimating which forecast
to trust adds more variance than the adaptivity recovers. \citet{hibon2005} posed
precisely our question in the pre-agentic setting, whether to select among forecasts
or combine them, and found selection hard to beat combination.

We state this at the outset because it sets the prior against which our central result
should be read. A reader who knows this literature should \emph{expect} a static
combination to beat an adaptive per-instance policy, and should regard a paper that
merely reports that direction as rediscovering a fifty-year-old result in new clothing.
Our claim is not the direction. It is that the agentic forecasting literature has been
building adaptive correction machinery without measuring it against that prior, that the
shortfall is large enough to matter, and --- the part classical work could not supply,
because it never had counterfactual outcomes for corrective \emph{actions} --- that we
can decompose the shortfall into named, separately quantified architectural causes and
say which one would have to be fixed first.

Alongside this, a separate literature has learned to \emph{predict} forecast failure.
Feature-based meta-learners estimate the error a model will incur on a given series
\citep{fformpp2019}, and stress-testing methods classify the conditions under which a
forecaster will produce unusually large errors, excessive uncertainty, or
``hubris''---large errors held with high confidence \citep{mast2025}. These methods stop
at diagnosis. Nobody has closed the loop by using a calibrated failure prediction to
decide whether an agent should act.

Uncertainty is likewise largely absent from agentic time-series systems. A recent survey
of reasoning and agentic systems for time series organises the field by reasoning
topology and objective, and contains no uncertainty category at all
\citep{tsreasoningsurvey2026}. Calibrated confidence and selective prediction are mature
elsewhere \citep{geifman2017selective,guo2017calibration} and well developed for time
series in the conformal literature \citep{gibbs2021aci,cptc2025}, but they have not been
connected to an agent's control decisions. The nearest work asks whether general-purpose
tool-using agents know when \emph{not} to act \citep{agentabstain2026}; it does not
address forecasting, and it has no counterfactual outcome data.

\subsection{What we find}

Four findings complicate the case for adaptive self-correction, and none of them is
what we set out to show. We lead with the one that bears on how correction policies are evaluated.

\paragraph{The oracle headroom that motivates correction is almost entirely
selection.}
Correction papers are usually judged by how much of an oracle's gap they close, where
the oracle picks the best action per instance on the test outcome. On our benchmark
that oracle promises a large improvement over accepting. Making the same choice on
validation origins and scoring it on later test origins recovers nothing: between
\rdCrpsCurseShare\% and \rdMaseCurseShare\% of the apparent gap is winner's
curse (Section~\ref{sec:decomposition}). Ground-truth failure labels on synthetic
data let us also grant a perfect diagnosis: applying the matched repair still does not
help, and on CRPS it hurts. The consequence is practical. A policy that closes half of
a test-selected gap may have closed none of the realizable one, so a correction method
should be measured against a validation-selected oracle and against a static default,
not against doing nothing.

This is a result about our ten-action menu and our hand-designed mapping from failure
mode to candidate repairs. It is not a test of any published system, none of which we
re-implemented.

\paragraph{Reflection does not restrain intervention.}
The systems this paper is about let a language model make the correction decision, so
we ran that too (Section~\ref{sec:llmarms}). On the primary model family every LLM
arm is worse than never correcting, with intervals excluding zero on both metrics; two
smaller families point the same way without resolving it. The reflection step fails by
being inert: asked whether a forecast needs revising, the model says yes on
\rateLlmCritiqueGate\% of instances, about as often as it intervenes with no
critique at all, and adding the critique changes how well it discriminates by
\discrimAurocDelta{}. At a matched rate of intervention, a small fitted estimator on the
same evidence decides better. The model's own stated risk is roughly two and a half
times the true rate of large errors, which is why it acts so often.

This is the forecasting counterpart of \citeauthor{huang2024selfcorrect}'s negative
result for reasoning, and it holds even though the evidence contained the external
signal, recent backtest error, that \citet{kamoi2024selfcorrect} identify as what
intrinsic self-correction lacks. The signal was there; the model did not use it.

\paragraph{The benefit is probabilistic, not point-accurate.}
The base forecasts are badly over-confident: nominal 80\% intervals achieve roughly 65\%
empirical coverage. Corrective actions repair this miscalibration -- and therefore CRPS
-- far more than they change the point forecast. A benchmark that scores point error
alone is close to blind to the effect.

\paragraph{A trivial static policy is not beaten by any adaptive policy we built.}
We designed reliability-gated correction to fix the known failure modes of naive
gating: a properly calibrated, risk-averse, per-action gain estimate, validated on
held-out backtest origins. It works better than a naively-thresholded gate. It does not
work better than the simplest possible alternative: applying a backtest-weighted
ensemble to \emph{every} instance, with no diagnosis, no gate, and no adaptivity at all.

All inference below is at the \emph{series} level --- one value per series, bootstrap
resampling over series --- as our pre-registration fixes, because instances drawn from
the same series share a generating process, a base model and most of their history, and
treating them as independent produces intervals that are far too narrow. At that level
the picture is sharp, and sharper than we expected:

\begin{itemize}
  \item On \textbf{point accuracy}, \emph{nothing works}. The static ensemble's
  \maseAlwaysEnsembleVsNeverPct\% improvement has a 95\% interval of
  [\maseAlwaysEnsembleVsNeverCIlo, \maseAlwaysEnsembleVsNeverCIhi]\%, which includes zero, as does
  our calibrated gate's. Meanwhile correcting indiscriminately is reliably
  \emph{harmful}: always-correct costs \maseAlwaysCorrectVsNeverPct\%
  ([\maseAlwaysCorrectVsNeverCIlo, \maseAlwaysCorrectVsNeverCIhi]\%) and a
  rate-matched random corrector \maseRandomCorrectVsNeverPct\%, both intervals excluding zero.
  \item On \textbf{probabilistic accuracy}, exactly one policy of any kind reliably
  helps, and it is the static one: the unconditional ensemble improves series-mean CRPS
  by \crpsAlwaysEnsembleVsNeverAbs\% ([\crpsAlwaysEnsembleVsNeverAbsCIlo,
  \crpsAlwaysEnsembleVsNeverAbsCIhi]\%). Our calibrated adaptive gate does \emph{not} clear the
  same bar: \crpsRgcCalibratedVsNeverPct\% with an interval of [\crpsRgcCalibratedVsNeverCIlo,
  \crpsRgcCalibratedVsNeverCIhi]\% that includes zero.
\end{itemize}

So the defensible claim is not that a static ensemble beats adaptive correction
everywhere. It is narrower and, we think, more useful: \emph{across both metrics, the
only reliable improvement available in this benchmark comes from a policy with no
adaptivity in it at all}, and every adaptive policy we built, including one designed
specifically to avoid the obvious failure modes, fails to establish an improvement
over doing nothing.

Two cautions belong with this. First, rank and magnitude disagree, as they do throughout
heavy-tailed forecasting data: the ensemble wins on \maseAlwaysEnsembleVsNeverWinRate\% of series on
MASE, which a signed-rank test detects easily, yet the sizes of those wins do not add up
to a mean improvement that excludes zero. We report both statistics everywhere. Second, as set out above, the \emph{direction}
here is what the combination literature predicts. What is new is the magnitude in the
agentic setting, the fact that this literature has not been measuring against that prior
at all, and the decomposition that follows.

These four findings are cautionary about \emph{automated} correction. Not every use of a
reliability signal needs to be. The same risk estimate that fails to drive good
correction decisions is, independently, a well-calibrated ranking signal for triage
(Section~\ref{sec:triage}): reviewing the riskiest tenth of forecasts by predicted risk
catches roughly a third of all large-forecast errors at two-thirds precision. We report
this as a positive, immediately deployable finding that does not depend on any of the
adaptive-correction machinery having worked.

\subsection{Contributions}

\begin{enumerate}
  \item \textbf{TS-AgentBench}, a benchmark carrying counterfactual ground truth for the
  correction decision. Every corrective action in a typed action space is executed on
  every instance and scored, yielding an oracle action, a \texttt{should\_correct}
  label, and a computable regret for any policy. To our knowledge no existing
  time-series benchmark supports decision-level evaluation of this kind. Because all
  outcomes are precomputed, evaluating a new correction policy costs a table lookup
  instead of a refit, which makes exhaustive ablation affordable and results exactly
  reproducible.

  \item \textbf{Reliability-gated correction (RGC)}, a policy in which the trigger is a
  calibrated quantitative estimate instead of a verbal judgement, the repair is
  conditioned on a diagnosed failure mode, and adoption requires demonstrated
  improvement on held-out backtest origins. This extends offline forecast-stress
  prediction \citep{fformpp2019,mast2025} into agent control, and grounds every step in
  an external signal as \citet{kamoi2024selfcorrect} indicate is necessary.

  \item \textbf{A compute-matched audit} of when self-correction helps, including a
  rate-matched random-correction control that spends an identical budget, per-regime
  analysis on data whose failure modes are known by construction, and reporting on both
  point and probabilistic accuracy. This audit shows that no adaptive policy we built
  beats never correcting, and that an unconditional static ensemble is the only policy
  of any kind that does, on probabilistic accuracy.

  \item \textbf{A demonstration that the oracle headroom motivating correction is
  almost entirely selection.} The oracle conventionally used to bound what correction
  could achieve picks the best action per instance on the test outcome. Choosing instead
  on validation origins and scoring on held-out test origins recovers nothing, and
  granting a perfect failure-mode diagnosis does not help either. Correction methods
  should therefore be measured against a validation-selected oracle and a static
  default; a policy that closes half of a test-selected gap may have closed none of the
  realizable one.

  \item \textbf{A positive, deployable finding alongside the cautionary ones}: the same
  failure-risk estimate that does not support good automated correction decisions is a
  well-calibrated triage signal, letting a practitioner trade a review budget for a
  quantified fraction of caught large-error forecasts.

  \item \textbf{Three measurement pathologies} that conceal the effect, each documented
  with the evidence that exposed it. First, expressing an action's value as an error
  \emph{ratio} divides by a baseline error that is sometimes near zero; on that scale the
  correlation between an action's observed and its realised value is $0.006$, while on a
  log scale the identical data gives $0.43$. Second, the arithmetic mean of MASE on this
  benchmark describes a few dozen extreme instances instead of the benchmark: the worst
  1\% of instances carry 19\% of the total error mass. Third, the oracle conventionally
  used to bound achievable headroom takes a minimum over many noisy candidate outcomes
  and is inflated by selection; the achievable, selection-free oracle is dramatically
  smaller. An evaluation of correction that ignores these will reach the wrong
  conclusion, as ours initially did.
\end{enumerate}

\subsection{What we do not claim}

The agentic forecasting area is crowded, and we are explicit about what is already
established. We do not claim the first agentic forecasting framework
\citep{tsscientist2025,nexus2026,castflow2026,tsagent2025,lastmile2026}; the first
time-series agent benchmark
\citep{temporalbench2026,timesagemt2026,timeseriesgym2025,aion2026}, nor the first
large-scale forecast evaluation suite \citep{gifteval2024,monash2021}; the first argument
that forecasting should be agentic \citep{position2026}; the first prediction of
forecast failure from features \citep{fformpp2019,mast2025}; the first evaluation of
agentic abstention \citep{agentabstain2026}; or, most importantly for reading our
headline result, the first demonstration that combining forecasts beats selecting
among them, which is fifty years old \citep{bates1969,clemen1989,wang2023review}. Our
contribution is narrower and, we
believe, currently missing: making the correction decision measurable, and testing
whether it survives a controlled comparison.

\subsection{Hypotheses}

The analysis was pre-registered before any experiment was run
(\texttt{docs/ANALYSIS\_PLAN.md} in the released code). We state the hypotheses here and
report every verdict in Section~\ref{sec:results}, including those not supported.

\begin{description}
  \item[H1] Reliability-gated correction improves accuracy relative to a non-correcting
  agent.
  \item[H2] Gains concentrate under distribution shift, anomalies and regime change.
  \item[H3] Forecast failure is predictable in advance from history-only diagnostics.
  \item[H4] Selective tool use matches exhaustive tool use at lower computational cost.
  \item[H5] Reliability-gated confidence is better calibrated than an agent's
  self-reported confidence.
  \item[H6] Any gain is not explained by additional computation alone.
  \item[H7] A quantitative gate decides better than verbal self-critique.
\end{description}

H6 and H7 are the load-bearing ones. Without H6, H1 is not falsifiable against the
``more compute'' explanation; we therefore fixed in advance that if H6 fails, H1 would
be reported as \emph{not attributable to the correction policy}, whatever H1 itself
showed.
